\section*{Preparation appendix: the arguments behind the notes}
This appendix is for rehearsal and questions. The timed talk is complete
at slide 27. The proof below uses the source notation only where useful;
in particular $M$ is never defined by dividing a character by two.

\subsection*{1. Source map and scope of the black boxes}
\begin{center}\small
\begin{tabular}{>{\raggedright\arraybackslash}p{.29\linewidth}>{\raggedright\arraybackslash}p{.35\linewidth}>{\raggedright\arraybackslash}p{.27\linewidth}}
Input & Exact source & Treatment\\\hline
Derivative frequencies and graph energy & \GT{Proposition 5.1} & Proved, slides 7--8\\
BSG and Pl\"unnecke & \GT{Theorems 5.2--5.3} & Stated black boxes, slide 9\\
Separation and graph refinement & \GT{Lemmas 8.3, 9.2} & Explained, slide 10\\
Global Bogolyubov & \GT{Lemma 6.3} & Stated, slide 11; proof in backup\\
Recovering correlation & \GT{Propositions 9.1, 9.3}; \JT{Lemma 2.5} & Proved, slides 11--12\\
Regular-radius existence & \GT{Lemma 8.2} & Stated black box, backup 1\\
Regularity shift estimates & \GT{Lemma 4.2} & Boundary proof in backup 1\\
Local Fourier decay & \GT{Lemma 8.4} & Proof, slide 15 notes / backup 4\\
Local Bogolyubov & \GT{Lemma 8.8} & Stated combinatorial input, slide 15; derivation below\\
Symmetry & \GT{Lemma 9.4}; \JT{Lemma 2.6} & Proved, slides 14--17\\
Lift and integrate & \JT{Lemmas 2.2, 2.4} & Stated algebraic black boxes, slides 21--23\\
Final Fourier coefficient & \GT{Lemma 4.4} & Proved, slide 26\\
Finite-field comparison & \Green{Propositions 2.5--2.6, Lemma 2.8} & Optional model, backup 2
\end{tabular}\end{center}
Here [JT] is arXiv:2112.13759v3 (2023); [GT] is arXiv:math/0503014v4
(December 2023); [Green] is arXiv:math/0604089v2 (2007). The bibliography
slide gives publication information and clickable links. The numbering in
this presentation is the main paper's numbering: the frequency-map result
is \emph{Lemma} 2.5, not Theorem 2.5. Green's notes have a different
Proposition 2.5.

\subsection*{2. The Cauchy--Schwarz calculation in slide 8}
Extend $\xi_h$ arbitrarily outside $H$ and put $F_k(h)=\mathbf1_H(h)e(\xi_h\cdot k)$.
For $a_h(x)=f(x+h)\overline{f(x)}$, squaring the Fourier coefficient gives
\[
|\widehat a_h(\xi_h)|^2
=\E_{x,k}f(x+h)\overline{f(x+h+k)}\overline{f(x)}f(x+k)e(\xi_h\cdot k).
\]
Thus, with $u(x,k)=f(x)\overline{f(x+k)}$ and
$v(x,k)=\overline{f(x)}f(x+k)$,
\[
\delta_0:=\E_h\mathbf1_H(h)|\widehat a_h(\xi_h)|^2
=\E_{x,h,k}u(x+h,k)v(x,k)F_k(h)\geq\eta^{16}/4.
\]
Since $|v|\leq1$, Cauchy--Schwarz gives
\[
\delta_0^2\leq
\E_{x,k}\abs{\E_hu(x+h,k)F_k(h)}^2
=\E_{y,a,k}u(y,k)\overline{u(y+a,k)}
       \E_hF_k(h)\overline{F_k(h+a)}.
\]
Here $y=x+h$ and $a$ is the difference of the two copies of $h$.
The right side is real and nonnegative before rearrangement. Since the
bounded factor is independent of $h$, a second Cauchy--Schwarz yields
\[
\delta_0^4\leq
\E_{a,k}\abs{\E_hF_k(h)\overline{F_k(h+a)}}^2
=\E_{h,a,b,k}F_k(h)\overline{F_k(h+a)}
                  \overline{F_k(h+b)}F_k(h+a+b).
\]
Average over $k$ using character orthogonality. The result is exactly the
normalised count of parallelograms in $H$ whose four frequencies obey the
same alternating relation. It is at least $\eta^{64}/256$. There are
$|G|^3$ ordered base parallelograms, proving the energy estimate with the
correct normalisation. This argument works for complex 1-bounded $f$ and
does not need the real-valued Samorodnitsky identity used in Green's notes.

\subsection*{3. Why the graph refinement has polynomial cost}
After BSG and Pl\"unnecke, enlarge $K\ll\eta^{-O(1)}$ until
$|\Gamma'|\geq |G|/K$ and $|9\Gamma'-8\Gamma'|\leq K|G|$.
Define
\[V=\{v\in\Gh:(0,v)\in8\Gamma'-8\Gamma'\}.\]
For each first coordinate of $\Gamma'$, adding all of $V$ produces
$|V|$ different points. Therefore
\[|\Gamma'||V|=|\Gamma'+(\{0\}\times V)|\leq K|G|,
\qquad |V|\leq K^2.\]
Characters on $\Gh$ are evaluations at elements of $G$, by finite
Pontryagin duality. If $v\ne0$, a uniformly random character $s$ of
$\Gh$ has $s(v)$ uniform on a nontrivial finite cyclic subgroup of
$\T$. At most two thirds of that subgroup lies in $(-1/4,1/4)$.
Successively choose a character leaving at most two thirds of the surviving
nonzero elements of $V$. After $O(\log(2|V|))$ choices,
\[V\cap\{v:\norm{s(v)}_\T<1/4\text{ for every chosen }s\}=\{0\}.\]
Let $m$ be the number of tests and $\Psi:\Gh\to\T^m$ their product.
Partition the torus into $128^m$ cubes of side $1/128$, using
representatives in $[0,1)^m$. Choose a cube containing at least this
fraction of the frequencies of $\Gamma'$ and call the resulting graph
$\Gamma''$. Since $m=O(\log K)$, its size is at least
$K^{-O(1)}|G|$.

If $(0,v)\in8\Gamma''-8\Gamma''$, then $v\in V$. In each tested
coordinate the eight positive and eight negative cube origins cancel;
the remaining real sum has magnitude below $16/128<1/4$.
Thus $v=0$. If two points of $4\Gamma''-4\Gamma''$ share a first
coordinate, their difference is such a vertical vector, so the two
points coincide. This proves injectivity of projection on that sumset.
We do \emph{not} claim that $8\Gamma''-8\Gamma''$ itself is a graph;
only its zero fibre is trivial.

\subsection*{4. From uniqueness of lifts to the full local-linearity condition}
Let $H''=\pi_G\Gamma''$ and choose $S$ by global Bogolyubov so that
$U_0=\Bohr(S,1/4)\subseteq2H''-2H''$. Define $M$ by the unique
lift $(h,M(h))\in2\Gamma''-2\Gamma''$ for $h\in U_0$.
For $h_1,h_2,h_3,h_4\in U_0$ with $h_1+h_2=h_3+h_4$, compare the
two sums of lifts in $4\Gamma''-4\Gamma''$. They have the same first
coordinate and hence the same second coordinate:
\[M(h_1)+M(h_2)=M(h_3)+M(h_4).\]
This is precisely the vanishing of all second differences. The
three-point additivity identity alone on an arbitrary subset would not
automatically imply this stronger four-point statement, so it is useful
to prove the latter directly. Also $M(0)=0$. The regular radius
$\rho\in[1/16,1/8]$ ensures $\Bohr(S,2\rho)\subseteq U_0$.

To get back to $f$, choose $x_0$ so the set
$A=\{h\in\Bohr(S,\rho):x_0+h\in H''\}$ has relative density
at least $|H''|/|G|$. Then $\xi_{x_0+h}=\xi_0+M(h)$ on $A$.
Write
\[c_h=\E_xf(x+x_0+h)\overline{f(x)}e(-[\xi_0+M(h)]\cdot x),\]
and choose $\lambda_h$ with $|\lambda_h|=1$ and
$\lambda_hc_h=|c_h|$ (choose $1$ if $c_h=0$). After $u=x+x_0$,
the explicit bounded factors can be taken as
\[
b_1(h)=\mathbf1_A(h)\lambda_h e((\xi_0+M(h))\cdot x_0),
\qquad b_2(u)=\overline{f(u-x_0)}e(-\xi_0\cdot u).
\]
This verifies both the size and the exact variable dependencies in
Lemma 2.5. No replacement of $f$ by an uncontrolled function is needed.

\subsection*{5. Symmetry: the algebra inside the Cauchy--Schwarz step}
After averaging translates, the Lemma 2.5 correlation gives
\[
\abs{\E_{h\in P,x\in U}a(h)c(x+h)b(x)e(-M(h)\cdot x)}\geq\delta,
\]
where $P$ is the original regular Bohr set, $U$ is much smaller and
regular, and $\delta\gg\eta^{O(1)}$. All points used below lie in the
original doubled domain of $M$. Squaring in $h$ introduces $x,y\in U$
and the phase $M(h)\cdot(y-x)$. Put $w=x+y+h$ and expand:
\[
\begin{aligned}
M(w-x-y)\cdot(y-x)
={}&M(w)\cdot y-M(w)\cdot x\\
&+M(x)\cdot x-M(y)\cdot y-D(x,y).
\end{aligned}
\]
Everything except $-D(x,y)$ is a sum of a function of $(w,x)$ and
a function of $(w,y)$. The factors $c(x+h)$ and
$\overline{c(y+h)}$ become $c(w-y)$ and $\overline{c(w-x)}$,
which have the same separated dependencies. Replace $P+x+y$ by $P$
using regularity, losing $O(d\,\operatorname{rad}(U)/\operatorname{rad}(P))$.
Fix $w$ by averaging and, if necessary, interchange $x,y$. Then
\[\abs{\E_{x,y\in U}u(x)v(y)e(D(x,y))}\geq\delta_1,
\qquad\delta_1\gg\eta^{O(1)}.\]
One more Cauchy--Schwarz gives
\[\E_{y,y'\in U}\abs{\E_{x\in U}e(D(x,y-y'))}\geq\delta_1^2.\]
There is a $y_0$ and a subset $A\subseteq U$ of relative density
$\gg\delta_1^2$ for which each inner absolute value with $y'=y_0$
is $\gg\delta_1^2$. Apply local Fourier decay to obtain
\[\norm{D(t,y_0-y)}_\T\ll\eta^{-O(1)}\norm{t}_S
\quad(t\in U,\ y\in A).\]
For $z=a_1+a_2-a_3-a_4\in2A-2A$, write
\[z=(y_0-a_3)+(y_0-a_4)-(y_0-a_1)-(y_0-a_2).\]
Bilinearity and the triangle inequality yield the same estimate with a
factor four. The ambient Bohr set was chosen sufficiently small that all
partial sums needed here lie in the original domain.

Local Bogolyubov now produces $S_*=S\cup T$ of polynomial size and a
radius $r_*\gg\eta^{O(1)}$ such that
$\Bohr(S_*,r_*)\subseteq U\cap(2A-2A)$. Consequently
\[\norm{D(t,z)}_\T\leq L\norm{t}_{S_*},\qquad L\ll\eta^{-O(1)},\]
for both variables in that Bohr set. Shrink to $2r\leq\varepsilon/L$
and $2r\leq r_*$, with $r$ regular and $\varepsilon$ a sufficiently
small power of $\eta$. This makes the real-midpoint argument valid on
the whole doubled Bohr set needed in Lemma 2.6.

\subsection*{6. Why local Bogolyubov has polynomial parameters}
This is the extra additive-combinatorial ingredient in the symmetry
step. Here is the mechanism of \GT{Lemma 8.8}, with constants suppressed
but the density dependence retained.

Let $U=\Bohr(S,r)$ be regular, $d=\max(1,|S|)$,
$A\subseteq U$, and $|A|=\delta|U|$. Take a regular
$U'=\Bohr(S,r')$ with $r'\asymp\delta r/d$, sufficiently small that
the regularity error is at most $\delta/2$. Averaging over translates
$x+U'$, $x\in U$, gives some $x_0$ for which
\[A'=A\cap(x_0+U'),\qquad |A'|\geq(\delta/2)|U'|.\]
Use $\beta=|U|/|G|$, $\beta'=|U'|/|G|$ and normalised convolution.
The function $\mathbf1_A*\mathbf1_{A'}$ is supported in
$x_0+U+U'$, of size at most $2|U|$. Cauchy--Schwarz and Parseval give
\[
F(0):=\sum_\xi|\widehat{\mathbf1_A}(\xi)|^2
                         |\widehat{\mathbf1_{A'}}(\xi)|^2
\geq\frac{\delta^4\beta(\beta')^2}{8}.
\]
The function
$F=\mathbf1_A*\mathbf1_{A'}*\mathbf1_{-A}*\mathbf1_{-A'}$
is supported in $2A-2A$. Define
\[R=\{\xi:|\widehat{\mathbf1_{A'}}(\xi)|\geq
                       (\delta^{3/2}/8)\beta'\}.\]
The Fourier mass outside $R$ is at most
$\delta^4\beta(\beta')^2/64$. Therefore, whenever
$\norm{\xi\cdot x}_\T<1/10$ for every $\xi\in R$,
\[
F(x)\geq\tfrac12F(0)-2\sum_{\xi\notin R}
|\widehat{\mathbf1_A}(\xi)|^2|\widehat{\mathbf1_{A'}}(\xi)|^2>0.
\]
It remains to enforce these possibly numerous frequency constraints with
only polynomially many new characters. That is the role of the local
Bessel inequality.

\textbf{Local Bessel input, exact claim.}
If $V=\Bohr(S,s)$ is regular, $E\subseteq V$, and
\[\mathcal R=\{\xi:|\widehat{\mathbf1_E}(\xi)|\geq t|V|/|G|\},\]
then, for $0<\theta,t\leq1$, there exist at most $2/t^2$ frequencies
$\xi_1,\ldots,\xi_m$ such that each $\xi\in\mathcal R$ satisfies
\[\sup_{x\in\Bohr(S,\theta s)}
\norm{(\xi-\xi_j)\cdot x}_\T\leq2^{16}\theta d/t^4
\quad\text{for some }j.\]
This is \GT{Corollary 8.7}. To see its proof, choose a maximal set of
frequencies separated by $q=2^{16}\theta d/t^4$ in this seminorm.
Local Fourier decay (applied to differences of characters) bounds the
off-diagonal inner products on $V$ by
$2^6(\theta d/q)^{1/2}\leq t^2/4$.
Choose phases aligning the correlations with $\mathbf1_E$.
Cauchy--Schwarz gives
\[t^2m^2\leq\E_{x\in V}\abs{\sum_{j=1}^m c_j e(\xi_j\cdot x)}^2
\leq m+\tfrac12t^2m^2.\]
Hence $m\leq2/t^2$. Maximality ensures that these frequencies cover
the large spectrum in the displayed seminorm. This derives the input
from the local Fourier-decay lemma, rather than invoking an additional
structural theorem.

Apply it to $E=A'-x_0\subseteq U'$; translation preserves Fourier
magnitudes. With $t=\delta^{3/2}/8$, there are
$m=O(\delta^{-3})$ representative frequencies. Set
$\theta=c\delta^6/d$ sufficiently small that the approximation error
is below $1/20$. On a Bohr set enforcing $\norm{\xi_j\cdot x}<1/20$
and the old constraints at radius $\theta r'$, every frequency in $R$
has circle norm below $1/10$. Thus the positive-convolution argument
puts this Bohr set inside $2A-2A$.

Since $r'\asymp\delta r/d$, this elementary tracking yields the slightly
weaker but fully sufficient radius $c\delta^7r/d^2$, with
$O(\delta^{-3})$ new characters. \GT{Lemma 8.8} states the sharper
$c\delta^6r/d$ radius quoted on slide 15. The talk may use that stated
black box; the derivation here proves a polynomial-radius version from
scratch, which is all the symmetry argument needs. In particular, no
global density such as $|A|/|G|$ enters the rank bound.

\subsection*{7. The midpoint is exactly bilinear, not approximately so}
On the small doubled domain, choose $d(y,z)\in(-\varepsilon,\varepsilon)$
lifting $D(y,z)$, with $\varepsilon<1/10$. Whenever $y_1,y_2,y_1+y_2,z$
are in the domain,
\[d(y_1+y_2,z)-d(y_1,z)-d(y_2,z)\in\Z\cap(-3\varepsilon,3\varepsilon)=\{0\}.\]
The same holds in $z$, and $d(z,y)=-d(y,z)$. Therefore
\[B(y,z)=M(y)\cdot z-\tfrac12d(y,z)\pmod1\]
is symmetric and locally bilinear. Furthermore
\[|e(-M(y)\cdot z)-e(-B(y,z))|
=|e(-d(y,z)/2)-1|\leq\pi\varepsilon.\]
Choose $\varepsilon$ below a small fixed fraction of the current
correlation lower bound. The phase replacement then loses at most half
that bound. The fact that $d$ is a small \emph{real} lift is indispensable:
there is no globally consistent half-map on $\T$.

\subsection*{8. Lifting and integration: check the local domains}
For $r<1/2$, each $x\in\Bohr(S,r)$ has a unique lift
$\widetilde x=(x,\theta_x)$ with $\norm{\theta_x}_\infty<r$.
If $x,y,x+y$ lie in this set and $3r<1$, the integer vector
$\theta_x+\theta_y-\theta_{x+y}$ has norm below one, hence vanishes.
Similarly, if all four vertices of a parallelogram lie in the set and
$4r<1$, their lifts satisfy the same parallelogram relation. On any cube
lying there, the lifted vertices consequently form an affine cube.
Restricting a global quadratic on $G_S$ therefore gives a locally
quadratic map in the full cube sense.

The sequence $0\to\Z^S\to G_S\to G\to0$ proves finite generation:
lift a finite generating set of $G$ and append the coordinate generators
of $\Z^S$. There may still be a finite torsion subgroup upstairs.
Lemma 2.4(i) applies without any torsion-free assumption.

For the cyclic example the cover is precisely $\Z$, and
$\widetilde B(n,m)=\alpha nm$ is real-valued. Reducing modulo one and
integrating gives $\widetilde\phi(n)=\alpha n^2/2\pmod1$.
For a genuinely nontrivial circle-valued local obstruction on $\Z/17\Z$,
one may take $\alpha=\sqrt2$, so that $17\alpha\notin\Z$.

\subsection*{9. The final phase bookkeeping and Fourier normalisation}
The expansion in slide 25 is worth checking once. Symmetry gives
\[
B(h+y+z,k-y)=B(h,k)-B(h,y)+B(y+z,k)-B(y,y)-B(y,z).
\]
Use $B(y,z)=\phi(y+z)-\phi(y)-\phi(z)$. With the minus sign in the
exponential, one can assign the factors explicitly:
\[
\begin{aligned}
a_{x,h,k}(y+z)&=b_1(x,h+y+z)e(-B(y+z,k)+\phi(y+z)),\\
b_{x,h,k}(y)&=b_2(x,k-y)e(B(h,y)+B(y,y)-\phi(y)-B(h,k)).
\end{aligned}
\]
The remaining factor is exactly $f(x+h+k+z)e(-\phi(z))$.
Thus no unwanted function depending only on $z$ is hidden in $b$.
After $x'=x+h+k$ and pigeonholing $h,k$, the desired factorisation holds.

For the final Fourier lemma let $p=|P|/|G|$, $q=|Q|/|G|$ and
$r=|P+Q|/|G|$. Extend
$F=g\mathbf1_Q$, $B=b\mathbf1_P$, $A=a\mathbf1_{P+Q}$ by zero.
Then, with the normalised transform convention,
\[
I:=\E_{y\in P,z\in Q}b(y)a(y+z)g(z)
=\frac1{pq}\sum_\xi\widehat F(\xi)\widehat B(\xi)\widehat A(-\xi).
\]
Cauchy--Schwarz and Parseval yield
\[
|I|\leq\frac{\sup_\xi|\widehat F(\xi)|}{pq}\sqrt{pr}
=\sqrt{r/p}\sup_\xi\abs{\E_{z\in Q}g(z)e(-\xi\cdot z)}.
\]
Since $r\geq p$, this implies the weaker factor $p/r$ stated in
\GT{Lemma 4.4}. Apply to $P=\bb_1$, $Q=\bb_2$ for each $x$.
If the previous joint average has magnitude $\geq\delta$, then
\[\E_x\sup_\xi\abs{\E_{z\in\bb_2}g_x(z)e(-\xi\cdot z)}
\geq\tfrac12\E_x|I_x|\geq\delta/2.\]
Choose a maximizer $\xi(x)$ (the dual group is finite). This avoids
an unnecessary extra pigeonhole step and proves the exact conclusion.

\subsection*{10. Reading cautions and quantitative bookkeeping}
The following are mathematical clarifications for preparation, not
digressions to deliver in the talk.
\begin{itemize}
\item In the printed \JT{Example 2.3}, an integer representative map
is described as a homomorphic section. For $N>1$, a homomorphism
$\Z/N\Z\to\Z$ is zero, so no such section exists. Use the unique
small representative locally, or an arbitrary set-theoretic section
globally. The example and the lifting argument then work exactly as
illustrated.
\item \GT{Propositions 9.1 and 9.3} write the graph as $(h,2M(h))$;
the implicit odd-order assumption is explicitly identified and removed
in \JT{Lemma 2.5, footnote 6}. We construct $(h,M(h))$ directly.
The final doubling manoeuvre in \GT{Lemma 9.4} is likewise omitted.
\item In \GT{Proposition 9.1}, the lift belongs to
$2\Gamma''-2\Gamma''$, not $2H''-2H''$; JT footnote 7 flags this
typo. The comparison of two pair sums requires the already proved
graph property of $4\Gamma''-4\Gamma''$, not a claim that
$8\Gamma''-8\Gamma''$ is a graph.
\item The last Fourier display in the printed JT proof uses the
$\rho_1$ label, although the extracted variable $z$ is averaged over
$\Bohr(S,\rho_2)$ in the preceding trilinear form. The application of
\GT{Lemma 4.4} gives the coefficient on the latter set, which satisfies
all the bounds required by the theorem. The slides use $\bb_2$
consistently.
\item The December 2023 correction on the
\href{https://arxiv.org/abs/math/0503014}{Green--Tao arXiv record}
concerns Proposition 3.2, a global extension claim. It is not used here.
The needed algebraic extension is the explicitly stated JT Lemma 2.2
on the lifted group.
\end{itemize}
Every rank enlargement and correlation loss in Lemmas 2.5--2.6 is
polynomial in $1/\eta$. Before integration, all radii have polynomial
lower bounds. The exponential shrink in Lemma 2.4(ii) is compatible
with the target radius because $|S|\ll\eta^{-O(1)}$.
Choose the later radii with fixed safety margins so all shifted
variables remain in doubled domains and all regularity errors are
smaller than the retained correlation. Products and compositions of
finitely many polynomial bounds are still polynomial; multiplying
the exponential lower radius by a polynomial does not change its
allowed form. This is the extent of constant bookkeeping appropriate
for a 45-minute talk.
